\documentclass[12pt]{article}

\usepackage[utf8]{inputenc}
\usepackage{amsmath, amssymb, amsthm}
\usepackage{geometry}
\geometry{a4paper, margin=2.5cm}

\title{Teorema de Descompunere Laurent}
\author{}
\date{}

\theoremstyle{definition}
\newtheorem{definition}{Definiție}

\theoremstyle{plain}
\newtheorem{theorem}{Teoremă}

\begin{document}

\maketitle

\begin{definition}
Fie $f$ o funcție olomorfă într-o vecinătate a punctului $z_0$. Spunem că $f$ are un
\emph{zero de ordin $k \ge 1$} în $z_0$ dacă


\[
f(z_0) = f'(z_0) = \cdots = f^{(k-1)}(z_0) = 0, 
\qquad f^{(k)}(z_0) \ne 0.
\]


Echivalent,


\[
f(z) = (z - z_0)^k g(z),
\]


unde $g$ este olomorfă în vecinătatea lui $z_0$ și $g(z_0) \ne 0$.
\end{definition}

\begin{definition}
Fie $f$ olomorfă în domeniul perforat $0 < |z - z_0| < r$. Spunem că $f$ admite o
\emph{serie Laurent} în jurul lui $z_0$ dacă există coeficienți $(a_n)_{n \in \mathbb{Z}}$ astfel încât


\[
f(z) = \sum_{n=-\infty}^{\infty} a_n (z - z_0)^n,
\]


iar seria converge pentru $0 < |z - z_0| < r$.
\end{definition}

\begin{theorem}[Descompunerea Laurent]
Fie $f$ o funcție olomorfă în domeniul perforat


\[
A = \{ z \in \mathbb{C} : 0 < |z - z_0| < R \}.
\]


Atunci $f$ admite o reprezentare unică sub forma unei serii Laurent:


\[
f(z) = \sum_{n=-\infty}^{\infty} a_n (z - z_0)^n,
\]


care converge pentru orice $z \in A$.

Coeficienții sunt dați de formula Cauchy:


\[
a_n = \frac{1}{2\pi i} \int_{\gamma} \frac{f(\zeta)}{(\zeta - z_0)^{n+1}} \, d\zeta,
\]


unde $\gamma$ este un cerc mic în jurul lui $z_0$ conținut în $A$.
\end{theorem}

\end{document}
